\section{Mitigation}\label{sec:mitigation}

As demonstrated in Section~\ref{sec:eval}, vulnerabilities due to shared memory 
double fetches are an insidious problem plaguing TAs, both for TAs currently 
deployed on COTS phones and future TAs written in Rust.
Analysis of the discovered 
vulnerabilities revealed that the affected TAs do not rely on zero-copy shared 
memory for their functionality. Instead, it appears the TA developer did not 
consider the \texttt{memref} buffer to be shared.

This is not the developer's fault. Instead, it is a failure of the GlobalPlatform standard and 
the TEE implementations, which obscure the fact that \texttt{memref} buffers may 
be shared. First, TEE implementations may or may not implement zero-copy shared 
memory. A TA developed for a TEE without zero-copy shared memory may later be 
deployed on a TEE with shared memory, or the TEE itself may switch to zero-copy 
shared memory in a later version. Furthermore, while there is a warning in the 
GlobalPlatform standard PDF~\cite{gpintapi}, the proposed solution is ineffective. The document 
suggests using \texttt{TEE\_CheckMemoryAccessRights} to check 
if a memory region is shared with a NWCA. By \emph{not} setting the \texttt{TEE\_MEMORY\_ACCESS\_ANY\_OWNER} 
flag, the API returns -1 if the memory is shared. Besides being convoluted, it 
is unclear to us whether TEEs support this flag, since we have not 
found a single TA in our dataset calling this API to check for shared memory. 

\begin{figure}[]
  \centering
  \centering
\begin{tikzpicture}%[every node/.style={font=\fontfamily{lmvtt}\selectfont\footnotesize}]
\small
%\node at (0.5,1.75) {\textbf{\stageone}};    

\coordinate (normal) at (-2,0);
\coordinate (secure) at (2,0);

%\coordinate (TAp) at (0,2);
\begin{scope}[shift=(secure)]
\coordinate (tap) at (0,1);
\coordinate (teep) at (0,-2);
\node[draw, thick, minimum width=3cm, minimum height=2.25cm] (ta) at (tap) {};
\node[draw, thick, minimum width=3cm, minimum height=3cm] (tee) at (teep) {};
\node[align=center] at ($(ta.north) + (-1.125cm,-0.25cm)$) {TA};
\node[align=center] at ($(tee.north) + (-1.125cm,-0.25cm)$) {TEE};

\begin{scope}[shift=(tap)]
\node[draw, thick, minimum width=2cm, minimum height=0.5cm, align=left] (register) at (0,0.2) {RegisterShm(cmd,\\bool shm[4])};
\node[draw, thick, minimum width=2cm, minimum height=0.5cm] (tainvoke) at (0,-0.7) {InvokeCommand};
\end{scope}

\begin{scope}[shift=(teep)]
\node[draw, thick, minimum width=2cm, minimum height=0.5cm] (cmd2shm) at (0,0.75) {cmd2shm};
\node[draw, thick, minimum width=2cm, minimum height=0.5cm] (isreg) at (0,-0.125) {is registered?};
\node[draw, thick, minimum width=2cm, minimum height=0.5cm] (copy) at (0,-1) {local copy};
\end{scope}

\end{scope}

\begin{scope}[shift=(normal)]
\node[draw, thick, minimum width=3cm, minimum height=2.35cm] (nwca) at (0,0.95) {};

\begin{scope}[shift=(nwca)]
\node[draw, thick, minimum width=2cm, minimum height=0.5cm, align=left] (nwcainvoke) at (0,-0.2) {InvokeCommand(\\  cmd=1,\\  params=[memref, .., \\.. , ..])};

\end{scope}

\node[align=center] at ($(nwca.north) + (-0.9cm,-0.25cm)$) {NWCA};

\draw[->, thick, dashed, >=stealth]
    ($(copy.east)$)
    -- ($(copy.east) + (0.9,0)$)
    -- ($(tainvoke.east) +(0.7,0.1)$)
    -- ($(tainvoke.east) +(0,0.1)$)
;

\draw[->, thick, dashed, >=stealth]
    ($(isreg.east)$)
    -- ($(isreg.east) + (0.7,0)$)
    node[pos=0.35, above, sloped] {yes}
    -- ($(tainvoke.east) +(0.5,-0.1)$)
    -- ($(tainvoke.east) +(0,-0.1)$)
;

\draw[->, thick, dashed, >=stealth]
    ($(nwcainvoke.south)$)
    -- ($(nwcainvoke.south) + (0,-2.05)$)
    -- ($(isreg.west) +(0,0)$)
    node[pos=0.4, above, sloped] {cmd, params}
;

\draw[->, thick, >=stealth]
    ($(register.west)$)
    -- ($(register.west) + (-0.4,0)$)
    -- ($(cmd2shm.west) +(-0.7,0)$)
    -- ($(cmd2shm.west) +(0,0)$)
;

\draw[<->, thick, >=stealth] ($(isreg.north) + (0,0)$) -- ($(cmd2shm.south) + (0,0)$) node[midway, above, sloped] {};
\draw[->, thick, >=stealth] ($(isreg.south) + (0,0)$) -- ($(copy.north) + (0,0)$) node[pos=0.5, right=0.05] {no};


\end{scope}
%\begin{scope}[shift=(snapshotsp)]
%    \node[draw, minimum width=1.5cm, minimum height=1.5cm, thick, fill=white] (eta3) at (-0.05,-0.1) {};
%    \node[draw, minimum width=1.5cm, minimum height=1.5cm, thick, fill=white] (eta3) at (0.0,-0.15) {};
%    \node[draw, minimum width=1.5cm, minimum height=1.5cm, thick, fill=white] (eta2) at (0.05,-0.2) {};
%    \node[draw, minimum width=1.5cm, minimum height=1.5cm, thick, fill=white] (eta1) at (0.1,-0.25) {};
%\end{scope}
%\node[minimum width=1.5cm, minimum height=0cm] (snapshotlabel) at ($(snapshotsp)+(-0.2,0.85)$) {snapshots};
%\node[draw, minimum width=1.5cm, minimum height=0cm, thick] (snappc) at ($(eta1)+(0,-0.25)$) {DF PC 2};
%
%%\node[draw, ellipse, thick, minimum width=1.75cm, minimum height=2.5cm] at (snapshotsp) {};
%%\draw (snapshotsp) ellipse (1cm and 1.5cm);
%\node[draw, minimum width=2cm, minimum height=2.75cm, thick] (emu) at (emuc) {};
%\begin{scope}[shift=(emuc)]
%    \node[draw, minimum width=1.5cm, minimum height=1.5cm, thick] (eta) at (0,-0.5) {};
%    \node[draw, minimum width=1.5cm, minimum height=0cm, thick] (eharness) at (0,0.65) {harness};
%\end{scope}
%\node[draw, minimum width=1.5cm, minimum height=0.5cm, thick] (efuzzer) at (efuzzerp) {fuzzer};
%\node[draw, minimum width=1.5cm, minimum height=1.5cm, thick] (eshm) at (eshmp) {};
%\node[draw, minimum width=1.5cm, minimum height=0cm, thick] (eshmaddr) at ($(eshmp)+(0,-0.15)$) {DF addr};
%\node[draw, minimum width=1.5cm, minimum height=0cm, thick] (edfpc1) at ($(eta)+(0,0.1)$) {DF PC 1};
%\node[draw, minimum width=1.5cm, minimum height=0cm, thick] (edfpc2) at ($(eta)+(0,-0.45)$) {DF PC 2};
%\node[align=center] at ($(emu.north) + (-0.3cm,-0.25cm)$) {emulator};
%\node[align=center] at ($(eta.north) + (-0.5cm,-0.25cm)$) {TA};
%\node[align=center] at ($(eshm.north) + (-0.3cm,-0.25cm)$) {SHM};
%\draw[->, thick, >=stealth] ($(efuzzer.west) + (0,0)$) -- ($(eharness.east) + (0,0)$) node[midway, above, sloped] {};
%\draw[->, thick, >=stealth] ($(eharness.south) + (0,0)$) -- ($(eta.north) + (0,0)$) node[midway, above, sloped] {};
%\draw[->, thick, dashed, >=stealth] ($(edfpc1.east) + (0,0)$) -- ($(eshmaddr.west) + (0,0.1)$) node[midway, above, sloped] {};
%\draw[->, thick, dashed, >=stealth] ($(edfpc2.east) + (0,0)$) -- ($(eshmaddr.west) + (0,-0.1)$) node[midway, above, sloped] {};
%\draw[thick] ($(eshmaddr.west) + (0,-0.1)$) circle (4pt);
%\draw[->, thick, >=stealth]
%    ($(emu.south)$)
%    -- ($(emu.south) + (0,-0.1)$)
%    -- ($(eta1.south) + (-0.1,-0.45)$)
%    node[pos=0.5, below, sloped] {merge \& save}
%    -- ($(eta1.south) +(-0.1,0)$);
%   

\end{tikzpicture}

  \caption{
    An overview of our mitigation. A TA may register \texttt{memref} buffers of 
    one or more specific commands as zero-copy shared memory using \texttt{RegisterShm}.
    When the NWCA invokes the TA, the TEE runtime checks if the input buffers were registered, 
    creating a local copy if not.
    }
  \label{fig:mitigation}
\end{figure}

Instead of leaving shared memory handling to developers and obscure APIs, we propose 
to mitigate this design weakness with an extension to the GlobalPlatform standard. This extension 
makes zero-copy shared memory explicit, removing ambiguity on shared memory and 
ensuring that double fetches on \texttt{memref} buffers are not a security issue by 
default. If a TA developer wishes to make use of zero-copy shared memory, the developer 
needs to explicitly register a \texttt{commandID},
\texttt{memref} buffer combination as zero-copy shared memory with the TEE runtime. 
For \texttt{memref} buffers not registered, the TEE runtime will create a local 
copy of the buffer and pass the local buffer to the TA. This extension changes 
the shared memory model to a fail-closed design, meaning that TAs which do not 
consider shared memory will \textbf{not} have double fetch vulnerabilities. 
Essentially, it makes the two scenarios outlined in Figure~\ref{fig:shmex} explicit 
to the developer. Our mitigation fixes all vulnerabilities discovered by SHMuzz, 
without changes to the existing TAs.

To demonstrate the feasibility of our mitigation, we implement it for OP-TEE, the 
reference GlobalPlatform-compliant TEE implementation. We add a new API \texttt{TEE\_RegisterShm}, 
which allows a TA to register an entry in the \texttt{params} array for a given 
\texttt{commandID} as shared. In its \texttt{TA\_CreateEntryPoint} function, the 
TA registers the \texttt{memref} that needs to be zero-copy shared memory. 
Listing~\ref{lst:mitigta} shows how a TA leverages this new API to request a
zero-copy shared memory buffer. When the TA's \texttt{TA\_InvokeCommandEntryPoint} function is called, the TEE user 
space runtime checks for each \texttt{memref} parameter, if for the given 
\texttt{commandID}, the \texttt{memref} was registered as zero-copy shared memory. 
If yes, the runtime passes the pointer to the shared memory to the TA. 
If not, the runtime makes a local copy of the 
buffer and sets the pointer in \texttt{memref} to the local buffer.
Figure~\ref{fig:mitigation} shows how our mitigation works.
Our patch to OP-TEE adds this mitigation in 140 lines 
of code.
We have submitted our mitigation proposal to the GlobalPlatform alliance and are 
in active discussion with them to update the GlobalPlatform TEE internal core 
specification to address shared memory.

\begin{figure}[h]
    \definecolor{pgreen}{rgb}{0,0.5,0}
\definecolor{pblue}{rgb}{0.13,0.13,1}
\definecolor{pred}{rgb}{0.9,0,0}
\definecolor{pgrey}{rgb}{0.46,0.45,0.48}
\definecolor{light-gray}{gray}{0.975}
\definecolor{darkgreen}{rgb}{0,0.5,0}
\definecolor{vscodepurple}{RGB}{128, 0, 128}
\definecolor{codepurple}{rgb}{0.58,0,0.82}

\lstset{escapechar=@,
                backgroundcolor = \color{light-gray},
                linebackgroundcolor={%
                \ifnum\value{lstnumber}>0
                    \ifnum\value{lstnumber}<30
                        \color{light-gray}
                    \fi
                \fi
                %\ifnum\value{lstnumber}>10
                %    \ifnum\value{lstnumber}<12
                %        \color{teal!40}
                %    \fi
                %\fi
                \ifnum\value{lstnumber}>15
                    \ifnum\value{lstnumber}<17
                        \color{blue!20}
                    \fi
                \fi 
                \ifnum\value{lstnumber}>20
                    \ifnum\value{lstnumber}<22
                        \color{red!20}
                    \fi
                \fi
                \ifnum\value{lstnumber}>22
                    \ifnum\value{lstnumber}<24
                        \color{red!20}
                    \fi
                \fi
                \ifnum\value{lstnumber}>23
                    \ifnum\value{lstnumber}<25
                        \color{red!20}
                    \fi
                \fi
                },
                tabsize=2,
                basicstyle=\fontencoding{T1}\selectfont\footnotesize,
                %basicstyle=\fontfamily{lmvtt}\selectfont\footnotesize,
                commentstyle=\color{darkgreen!80},
                keywordstyle=\color{pblue},
                stringstyle=\color{pred},
                %stringstyle=\color{vscodepurple},
                frame = single,
                showstringspaces=false,
                language=c,
                xleftmargin=5pt, 
                xrightmargin=5pt, 
                %framexleftmargin=-5pt,
                numbers=left,                  % Line numbers
                numberstyle=\tiny\color{black}, % Line number style
                numbersep=4pt, %moredelim=[is][\textbf]{@}{@}
                %escapeinside={(*@}{@*)}
}
\center

\begin{lstlisting}[columns=fullflexible]
TEE_Result TA_CreateEntryPoint()
{
   bool paramsShm[4] = {true,false,false,false};
   TEE_RegisterShm(3, &paramsShm);
}
\end{lstlisting}

    \captionof{lstlisting}{
        The updated \texttt{TA\_CreateEntryPoint} function of the TA from 
        Listing~\ref{lst:exdf}, which explicitly registers the first \texttt{memref}
        buffer when handling command 3 as zero-copy shared memory. All the other 
        \texttt{memref} buffers are not zero-copy, and thus the double fetch bugs 
        in commands 1 and 2 are mitigated.
    }
    \label{lst:mitigta}
\end{figure}